% !TeX root = ../papers/nc-p3.tex
\section{階数15の座標端で同型の重複が残る有限水準}
\label{sec:rank-fifteen-coordinate-collisions}
\begin{theorem}\label{thm:rank-fifteen-coordinate-collisions}
前節の同時反転で割った三曲線を $E,R_+,R_-$ とし、各点の水準を $U=zk^2$ とする。
二つの相異なる通常な階数15の点が添字保持同型を与えるなら、
表の該当欄の多項式のいずれかが零になる。
二つの有理成分は互いに共役である。最後の行では添字で水準を区別する。
\[
\begin{array}{c|l}
\text{比較する曲線} & \text{水準の必要条件を与える因子} \\\hline
 E,E & U+3,\ U-3,\ U-9,\ U^2-74U+169\\
 R_+,R_+\ \text{又は}\ R_-,R_- & U+3,\ U+1,\ U^2-42U+9\\
 R_+,R_- & U+3,\ U-9\\
 E,R_\pm &
 U_E^3+13U_E^2+259U_E-529,\quad U_R^3-147U_R^2-61U_R-49
\end{array}
\]
最後の行では二つの三次式がともに消えなければならない。
従ってこの表の外では、三つの座標端を合わせても、前節の商パラメータが添字保持同型類を
一意に定める。表の零点における同型の存在はこの必要条件からは従わない。
\end{theorem}
\begin{proof}
$E$ の座標を $s=t+1/t$, $v^2=(14-s)(s^2-4)$ とし、有理曲線では $s=t-1/t$ を使う。
$R_+$ では $i^2=-1$ を一つ固定し、$R_-$ では $i$ を $-i$ に替える。水準は
\[
 U_E=\frac{(s+10)^2}{16(s+2)},\qquad
 U_{R_+}=\frac{i(s^2+4)}{8s}.
\]
通常域ではこれらの分母は零にならない。
節\ref{sec:rank-fifteen-endpoint-orbits}の四不変量 $J_2,J_3,C,L$ を各曲線の関数体で求める。
\path{computations/verify_rank_fifteen_coordinate_collisions.py} は五重縮約の
各文字の和が零になる条件を、置換 $Q$ の二閉路に沿って解く。
こうして全成分の積を整数多項式で足し、最後に共通分母を戻す。
結果を $t^2-st\pm1$ で簡約すると、$J_2,J_3,L$ は全ての曲線上で $s$ の有理関数であり、
$E$ 上の $C$ は $v$ の一次式になる。その $v$ の係数は
\[
 \frac{-48(s+2)(s^2-28s+4)}{(s^2+68s+196)^3}.
\]
従って通常域で $s$ が一致して $v$ が異なる二点は、トレースが非零なら $U_E=3$ でしか
同型になり得ない。

二点の $s$ 座標を $s,u$ とする。各不変量の両曲線上の表示を既約分数 $A/B,A'/B'$ とし、
等式の分子 $A(s)B'(u)-B(s)A'(u)$ を取る。
$J_2,J_3,L$ を常に使い、双方が有理曲線なら $C$ も使う。同一成分の比較では
各行を $s-u$ で割る。得た行の全ての対について $u$ を消去し、非零な終結式の
最大公約式の平方自由部分を取る。非通常点、階数14の $U=1$ と
$\operatorname{tr}S=0$ の因子を除いた後、水準式で $s$ を消去する。
コードは五方向の比較で、表の因子のうちこの消去から来るものを最高係数1で検算する。
共通零点は各非零終結式を消すので、この操作は必要条件を失わない。

残る $\operatorname{tr}S=0$ は三成分とも $U=-3$ に一致する。
ここで $\operatorname{tr}(S^2)$ は非零であり、
$\operatorname{tr}(S^3)^2/\operatorname{tr}(S^2)^3$ の値は $E$ で $64/375$、
両有理成分で $-64/81$ となる。これは $J_3^2/J_2^3$ を約分して $U=-3$ の方程式で
剰余を取ると得られる。従ってこの水準でも楕円と有理の間の同型は除かれる。
\end{proof}
